\documentclass[a4paper,10pt]{article}
\usepackage[margin=2.5cm]{geometry}
\usepackage[utf8]{inputenc}
\usepackage[colorlinks=true,urlcolor=blue]{hyperref}
\usepackage{amsmath}
\usepackage{graphicx}
\usepackage{float}
\usepackage{caption}
\usepackage{subcaption}
\usepackage{minted}



\setlength{\parindent}{0em}
\setlength{\parskip}{1em}

\begin{document}
\begin{titlepage}
    \centering
    \vspace*{2cm}
    {\Huge \textbf{1RT730 Project Report} \par}
    \vspace{0.5cm}
    {\LARGE Harald - Your Personal Chat-Powered Language Partner \par}
    
    \vspace{2cm}
    \includegraphics[width=0.6\textwidth]{cute_thinking_harald.png} % your image
    \vspace{2cm}
    
    {\Large Joel Sundin, Petter Möllerström, Rahul Sebastian Peter \par}
    \vspace{1cm}
    {\large \today \par}
\end{titlepage}   
\newpage

 \section{Introduction}
Learning Swedish as a non-native speaker presents significant challenges, particularly when opportunities for interactive practice and real-time feedback are limited. To address these challenges, this work proposes \textit{Harald}, an AI-driven language learning application that leverages Large Language Models (LLMs), quizzes, and flashcard functionalities to enhance engagement and personalization in language acquisition.

Harald integrates multiple learning modalities to support comprehension, vocabulary retention, and conversational skills. Beyond individual learning, Harald contributes to societal language inclusivity by supporting immigrants, students, and professionals in their integration into Swedish society. The primary contribution of this study lies in the integration of LLM-based interaction within a structured learning framework, enabling personalized practice, targeted feedback, and contextual understanding in a scalable and accessible manner.

Harald is implemented as a web-based application, allowing learners to access its features directly through a standard browser without the need for local installation. This design choice enhances accessibility and convenience, enabling learners to practice Swedish anytime and anywhere, whether on desktop or mobile devices. By combining cloud-based LLM computation with an intuitive browser interface, Harald ensures that advanced language processing capabilities are delivered seamlessly to users, while maintaining a lightweight and user-friendly experience.

The system’s web-based nature also facilitates easy updates and integration of new exercises or models, ensuring that learners always benefit from the latest improvements in AI-driven language education. Overall, Harald provides a comprehensive, adaptive, and accessible environment for Swedish language learning, bridging the gap between traditional pedagogical methods and modern AI-assisted instruction.


\section{Dataset}

The development of the Swedish language tutoring chatbot involved several iterations of dataset design and experimentation aimed at improving the model’s performance, naturalness, and pedagogical value. Over the course of the project, multiple datasets were constructed and evaluated, ranging from community-sourced linguistic discussions to synthetically generated dialogue data. Ultimately, however, the final system adopted a different strategy that relied on prompt engineering and curated conversational examples rather than direct fine-tuning. The main dataset iterations are described below.

\subsection{Reddit-based corpus}
The initial dataset was collected from the r/swedish subreddit using the Python Reddit API Wrapper (PRAW). This community consists of users interested in learning or discussing the Swedish language, and thus provides an organic dataset composed of learner questions and native or advanced speaker responses. The dataset comprised approximately 6,500 comments, each representing short-form discussions about grammar, vocabulary, pronunciation, and translation. Preliminary fine-tuning experiments were conducted using this dataset to adapt a general-purpose language model to the linguistic style and content of Swedish language learners. However, despite its relevance, the dataset exhibited several limitations, including inconsistent quality, informal tone, and limited conversational depth. Consequently, it was deemed insufficient as a primary fine-tuning source.

\subsection{Synthetic dialogue dataset}
In the second iteration, a manually constructed dataset of user–assistant interactions was developed to simulate tutoring dialogues. These pairs were designed to represent common learner scenarios such as vocabulary practice, grammar correction, and cultural explanations. This dataset was later expanded using generative augmentation, where additional examples were synthesized with Anthropic’s Claude Sonnet model. The resulting dataset provided higher-quality, domain-specific examples and was used for initial LoRA-based fine-tuning experiments. While the synthetic data improved linguistic consistency, further evaluation revealed limited benefits relative to the computational cost of model adaptation.

\subsection{Web-scraped Swedish news dataset}
To provide learners with authentic, up-to-date material for reading comprehension and writing practice, a web scraper was implemented to collect recent Swedish news articles from online media outlets. The system allows users to choose from a selection of articles and engage in follow-up exercises such as summarization or open-ended question answering. This component operates dynamically at runtime rather than as part of the fine-tuning dataset, ensuring continuous exposure to contemporary language usage and diverse topics.

\subsection{Flashcard-based vocabulary dataset}
An auxiliary feature enables users to store specific words encountered during interaction into a personal flashcard list for later review. Although not a dataset in the traditional training sense, this mechanism contributes to personalized learning and vocabulary retention through spaced repetition.

\subsection{Final data strategy}
In the final system, neither the Reddit corpus nor the synthetic fine-tuning datasets were employed in the deployed model. Instead, the chatbot is based on a fully pretrained large language model (Gemini 2.5 Flash), with its pedagogical behavior governed by an extensive system prompt. This prompt defines the tutor’s persona, instructional goals, and tone, supplemented by curated example dialogues that demonstrate effective interaction patterns. This approach provided greater flexibility, reduced training costs, and preserved the model’s general linguistic competence while enabling targeted guidance for Swedish language tutoring.

\section{Architecture}
This section should outline the overall design and workflow of the LLM-based application. Present the system architecture, describing how inputs are processed, how the LLM is integrated, and how outputs are generated and refined. If applicable, explain supporting components such as retrieval-augmented generation (RAG) modules, safety filters, knowledge bases, or user interfaces. Visual aids such as figures or flowcharts are encouraged to illustrate the data flow and interactions between components. Also please justify the design choices and explain how the architecture supports the application's goals.
\subsection{Previous Attempts}

\section{Results}
In this section, present the outcomes of your experiments or evaluations. Provide both quantitative results (e.g., accuracy, precision, recall, BLEU, latency) and qualitative analyses (e.g., case studies, user feedback, example outputs). Where possible, compare the performance of your system against baseline methods or alternative approaches. Discuss the strengths of your application as well as its limitations, including cases where it may fail or underperform. Use tables, graphs, or figures to support clarity and emphasize key findings.

\section{Societal Impact}
The use of large language models (such as Gemini 2.5 Flash) has a strong influence on society. These models can bring many benefits but also create new challenges that must be handled responsibly.

\subsection{Positive Implications}
The main benefit of Harald(Swedish chatbot) is improved accessibility and inclusion. Learners from different backgrounds can practice Swedish anytime without the need for expensive language courses or native-speaking tutors. This can support immigrants, students and professionals who want to integrate into Swedish society more easily. The chatbot also improves educational efficiency offering instant corrections, feedback and explanations in simple English. It encourages personalized learning allowing users to progress at their own pace and supports the democratization of knowledge by making language learning resources available to a wider audience.

\subsection{Negative Implications}
However there are several potential risks to be addressed. Since Harald uses the Gemini model, it relies on a pre-trained dataset that we do not have any control over. This means potential biases, inaccuracies or cultural insensitivities from the training data can surface in responses. The only adjustable parameter is the system prompt, which helps guide the model’s tone and focus but it might not be effective in all edge cases. There are also parameters which restricts harmful and hateful dialogues but inspite of all these techniques it is difficult to guarantee perfect accuracy or fairness in all interactions. There are also privacy concerns because users might share personal details in their messages. This requires careful data handling and protection. Another risk is that users could become too dependent on the chatbot or overconfident on AI-generated responses and use it instead of real conversations with people. In addition using large models like Gemini has an environmental cost, since it requires high energy use for both training and running the model.

\subsection{Ethical Considerations and Mitigation}
To manage these issues the system should be transparent about being an AI and not a real teacher. Users should understand that it may not always be correct or unbiased. The system prompts should be regularly updated to make the responses more reliable and human reviews can help check for mistakes or bias. Personal information should be anonymized or removed and data storage should follow privacy regulations. Using smaller or optimized versions of the model can also help lower the environmental impact.

While Harald highlights the educational and social benefits of integrating LLMs into language learning, ongoing monitoring is essential. Balancing accessibility and personalization with fairness, accuracy and ethical responsibility will determine the long-term success and trustworthiness of such systems in society.
	
\hfill \break
\textit{\textbf{Use of generative AI}}

Generative AI was used for formatting the text for writing this report. It was also used to create examples chatbot conversation to use as system prompts for making the chatbot accessible and personlized,


\hfill \break


\begin{thebibliography}{1}

\bibitem{google_gemini}
Google AI. \textit{Gemini API Text Generation Documentation.} 
Available at: \url{https://ai.google.dev/gemini-api/docs/text-generation}

\end{thebibliography}


	
	
%	\pagebreak
%	
%	\appendix
%	\section{Code}
%	\label{app:code}
%	
%	\inputminted[frame=lines,framesep=2mm,linenos]{python}{your_code.py}
%	
\end{document}
